\markboth{Abstract}{}
\noindent{\sffamily\fontsize{38}{44}\selectfont\bfseries\color{navy}EAR}\hfill{\sffamily\footnotesize\color{muted}Technical Report\quad /\quad 2026.10}\par
\vspace{1mm}\noindent{\sffamily\small\color{muted}Effective Auto Research}\par
\vspace{4mm}\noindent{\sffamily\fontsize{20}{28}\selectfont\bfseries\color{ink}Advance research through feedback.\\Improve strategies through results.}\par
\vspace{4mm}\noindent\AuthorsCompact
\vspace{3mm}\noindent\textcolor{navy!35}{\rule{\textwidth}{.5pt}}\par
\vspace{1mm}\subsection*{Abstract}
EAR (Effective Auto Research) is an open-source research automation system for idea discovery, autonomous mathematical research, and author responses to peer review. It aims to reduce researchers' active involvement while meeting research quality requirements within a machine budget. Research often advances when plans change: a new paper alters the choice of problem, a counterexample redirects a proof, or a review sends a manuscript back for revision. Existing systems have brought automation to ideation, experimentation, and writing. Sustained research also requires materials to carry forward, feedback to guide the next action, and experience from one task to inform the next.

EAR follows one principle: \textbf{use feedback to advance research, and research outcomes to improve its strategies}. AutoVibeIdea combines literature retrieval, critique, and tree search to turn a research direction into an executable proposal. AutonomousMath uses an inner loop for derivation, validation, and manuscript revision, and an outer loop to test and evolve execution strategies through complete research episodes. AutoRebuttal starts from reviewer concerns and available evidence, fills gaps, and compares and revises responses. The three modules can be used independently; researchers can also carry proposals, manuscripts, and review materials between stages. People retain control of direction, key technical checks, and final review.\src{1}

In historical evaluations, sustained revision increased current-version model-assessment passes from 3 to 12 among the same 28 manuscripts. A selected evolved strategy produced more passes than the initial strategy while requiring 52.8\% fewer research invocations per pass. These findings provide initial support for EAR's approach to sustained research and strategy improvement.\src{10}

\vspace{3mm}
\begingroup\small\sffamily\raggedright
\noindent\textbf{Keywords}\quad Research automation; active human time; critical tree search; strategy evolution; evidence-based rebuttal\par
\endgroup
\clearpage
\markboth{Contents}{}
\tableofcontents
\thispagestyle{fancy}
\vspace{8mm}
\begin{scope}
\textbf{Reading guide}\quad Section~\ref{sec:introduction} introduces EAR's scope and central principle. Section~\ref{sec:methods} explains the three workflows, with particular attention to the mathematical research loops. Section~\ref{sec:evaluation} presents manuscript revision, strategy comparisons, and rating-change prediction, together with their evaluation conditions.
\end{scope}
\clearpage
